% =====================================================================
%  Anatomía de un LLM — Fase II — Módulo 8
%  Aprendizaje supervisado
%  Compilar: latexmk -r ../latexmkrc M08-supervisado.tex
% =====================================================================
\documentclass[11pt,a4paper,twoside,openany]{book}
\usepackage{../estilo-apuntes}

\newcommand{\ApuntesCurso}{Anatomía de un LLM}
\newcommand{\ApuntesModulo}{Módulo 8}

% Anclas de hiperenlace únicas por tema para los enunciados
\makeatletter
\@for\nombre:=definicion,teorema,proposicion,lema,corolario,ejemplo,observacion\do{%
  \expandafter\xdef\csname theH\nombre\endcsname{\noexpand\thechapter.\noexpand\arabic{definicion}}}
\makeatother

% Notación propia del módulo
\DeclareMathOperator*{\argmin}{arg\,min}
\DeclareMathOperator*{\argmax}{arg\,max}
\DeclareMathOperator{\Var}{Var}
\DeclareMathOperator{\Cov}{Cov}
\DeclareMathOperator{\Ber}{Ber}
\DeclareMathOperator{\Cat}{Cat}
\DeclareMathOperator{\diag}{diag}
\DeclareMathOperator{\tr}{tr}
\DeclareMathOperator{\rg}{rg}
\DeclareMathOperator{\sign}{sign}
\DeclareMathOperator{\SSE}{SSE}
\DeclareMathOperator{\ECM}{ECM}
\DeclareMathOperator{\BCE}{BCE}
\newcommand{\N}{\mathcal{N}}
\newcommand{\D}{\mathcal{D}}
\newcommand{\Hc}{\mathcal{H}}
\newcommand{\X}{\mathcal{X}}
\newcommand{\Y}{\mathcal{Y}}
\newcommand{\ind}{\mathbb{I}}
\newcommand{\Remp}{\widehat{R}_n}

\hypersetup{pdftitle={Módulo 8. Aprendizaje supervisado (resumen)},pdfauthor={Anatomía de un LLM}}

\makeatletter
\newlabel{ej:bayes}{{1.8}{2}{}{ejemplo.1.8}{}}
\newlabel{ej:bayes@cref}{{[ejemplo][8][]1.8}{[1][2][]2}}
\newlabel{fig:particion}{{1.1}{4}{Protocolo de entrenamiento, validación y prueba. Las proporciones son las de la práctica del módulo}{figure.1.1}{}}
\newlabel{fig:particion@cref}{{[figure][1][1]1.1}{[1][3][]4}}
\newlabel{eje:1-tipos}{{1.1}{4}{}{ejercicio.1.1}{}}
\newlabel{eje:1-tipos@cref}{{[ejercicio][1][1]1.1}{[1][4][]4}}
\newlabel{eje:1-particion}{{1.2}{4}{}{ejercicio.1.2}{}}
\newlabel{eje:1-particion@cref}{{[ejercicio][2][1]1.2}{[1][4][]4}}
\newlabel{eje:1-bayes}{{1.3}{4}{}{ejercicio.1.3}{}}
\newlabel{eje:1-bayes@cref}{{[ejercicio][3][1]1.3}{[1][4][]4}}
\newlabel{eje:1-constante}{{1.4}{4}{}{ejercicio.1.4}{}}
\newlabel{eje:1-constante@cref}{{[ejercicio][4][1]1.4}{[1][4][]4}}
\newlabel{eje:1-logloss}{{1.5}{4}{}{ejercicio.1.5}{}}
\newlabel{eje:1-logloss@cref}{{[ejercicio][5][1]1.5}{[1][4][]4}}
\newlabel{ej:recta}{{2.8}{6}{Ajuste de una recta}{ejemplo.2.8}{}}
\newlabel{ej:recta@cref}{{[ejemplo][8][]2.8}{[1][6][]6}}
\newlabel{fig:proyeccion}{{2.1}{7}{Mínimos cuadrados como proyección ortogonal: $\hat y$ es el punto del espacio de columnas más próximo a $y$, y el residuo $r$ es perpendicular a ese espacio}{figure.2.1}{}}
\newlabel{fig:proyeccion@cref}{{[figure][1][2]2.1}{[1][7][]7}}
\newlabel{eje:2-rango}{{2.1}{9}{}{ejercicio.2.1}{}}
\newlabel{eje:2-rango@cref}{{[ejercicio][1][2]2.1}{[1][9][]9}}
\newlabel{eje:2-recta}{{2.2}{9}{}{ejercicio.2.2}{}}
\newlabel{eje:2-recta@cref}{{[ejercicio][2][2]2.2}{[1][9][]9}}
\newlabel{eje:2-strang1}{{2.3}{9}{}{ejercicio.2.3}{}}
\newlabel{eje:2-strang1@cref}{{[ejercicio][3][2]2.3}{[1][9][]9}}
\newlabel{eje:2-strang2}{{2.4}{9}{}{ejercicio.2.4}{}}
\newlabel{eje:2-strang2@cref}{{[ejercicio][4][2]2.4}{[1][9][]9}}
\newlabel{eje:2-laplace}{{2.5}{9}{}{ejercicio.2.5}{}}
\newlabel{eje:2-laplace@cref}{{[ejercicio][5][2]2.5}{[1][9][]9}}
\newlabel{eje:2-traza}{{2.6}{9}{}{ejercicio.2.6}{}}
\newlabel{eje:2-traza@cref}{{[ejercicio][6][2]2.6}{[1][9][]9}}
\newlabel{ej:logistica}{{3.7}{11}{}{ejemplo.3.7}{}}
\newlabel{ej:logistica@cref}{{[ejemplo][7][]3.7}{[1][11][]11}}
\newlabel{eje:3-lineal}{{3.1}{13}{}{ejercicio.3.1}{}}
\newlabel{eje:3-lineal@cref}{{[ejercicio][1][3]3.1}{[1][12][]13}}
\newlabel{eje:3-valores}{{3.2}{13}{}{ejercicio.3.2}{}}
\newlabel{eje:3-valores@cref}{{[ejercicio][2][3]3.2}{[1][13][]13}}
\newlabel{eje:3-gradiente}{{3.3}{13}{}{ejercicio.3.3}{}}
\newlabel{eje:3-gradiente@cref}{{[ejercicio][3][3]3.3}{[1][13][]13}}
\newlabel{eje:3-frontera}{{3.4}{13}{}{ejercicio.3.4}{}}
\newlabel{eje:3-frontera@cref}{{[ejercicio][4][3]3.4}{[1][13][]13}}
\newlabel{eje:3-newton}{{3.5}{13}{}{ejercicio.3.5}{}}
\newlabel{eje:3-newton@cref}{{[ejercicio][5][3]3.5}{[1][13][]13}}
\newlabel{eje:3-separable}{{3.6}{13}{}{ejercicio.3.6}{}}
\newlabel{eje:3-separable@cref}{{[ejercicio][6][3]3.6}{[1][13][]13}}
\newlabel{ej:softmax}{{4.4}{15}{}{ejemplo.4.4}{}}
\newlabel{ej:softmax@cref}{{[ejemplo][4][]4.4}{[1][15][]15}}
\newlabel{fig:cabeza}{{4.1}{16}{La capa de salida de un LLM es una regresión softmax sobre el vocabulario. En entrenamiento se evalúa la entropía cruzada del token real; en generación se muestrea de $p$ (módulo~1)}{figure.4.1}{}}
\newlabel{fig:cabeza@cref}{{[figure][1][4]4.1}{[1][16][]16}}
\newlabel{eje:4-llm}{{4.1}{17}{}{ejercicio.4.1}{}}
\newlabel{eje:4-llm@cref}{{[ejercicio][1][4]4.1}{[1][17][]17}}
\newlabel{eje:4-logits}{{4.2}{17}{}{ejercicio.4.2}{}}
\newlabel{eje:4-logits@cref}{{[ejercicio][2][4]4.2}{[1][17][]17}}
\newlabel{eje:4-k2}{{4.3}{17}{}{ejercicio.4.3}{}}
\newlabel{eje:4-k2@cref}{{[ejercicio][3][4]4.3}{[1][17][]17}}
\newlabel{eje:4-llama}{{4.4}{17}{}{ejercicio.4.4}{}}
\newlabel{eje:4-llama@cref}{{[ejercicio][4][4]4.4}{[1][17][]17}}
\newlabel{eje:4-gradiente}{{4.5}{18}{}{ejercicio.4.5}{}}
\newlabel{eje:4-gradiente@cref}{{[ejercicio][5][4]4.5}{[1][17][]18}}
\newlabel{eje:4-hessiano}{{4.6}{18}{}{ejercicio.4.6}{}}
\newlabel{eje:4-hessiano@cref}{{[ejercicio][6][4]4.6}{[1][18][]18}}
\newlabel{ej:encoger}{{5.6}{20}{Encoger reduce el error}{ejemplo.5.6}{}}
\newlabel{ej:encoger@cref}{{[ejemplo][6][]5.6}{[1][20][]20}}
\newlabel{eje:5-vf}{{5.1}{22}{}{ejercicio.5.1}{}}
\newlabel{eje:5-vf@cref}{{[ejercicio][1][5]5.1}{[1][22][]22}}
\newlabel{eje:5-diagnostico}{{5.2}{22}{}{ejercicio.5.2}{}}
\newlabel{eje:5-diagnostico@cref}{{[ejercicio][2][5]5.2}{[1][22][]22}}
\newlabel{eje:5-mano}{{5.3}{22}{}{ejercicio.5.3}{}}
\newlabel{eje:5-mano@cref}{{[ejercicio][3][5]5.3}{[1][22][]22}}
\newlabel{eje:5-interpolar}{{5.4}{22}{}{ejercicio.5.4}{}}
\newlabel{eje:5-interpolar@cref}{{[ejercicio][4][5]5.4}{[1][22][]22}}
\newlabel{eje:5-encoger}{{5.5}{22}{}{ejercicio.5.5}{}}
\newlabel{eje:5-encoger@cref}{{[ejercicio][5][5]5.5}{[1][22][]22}}
\newlabel{eje:5-varianza}{{5.6}{22}{}{ejercicio.5.6}{}}
\newlabel{eje:5-varianza@cref}{{[ejercicio][6][5]5.6}{[1][22][]22}}
\newlabel{ej:ridge}{{6.4}{24}{}{ejemplo.6.4}{}}
\newlabel{ej:ridge@cref}{{[ejemplo][4][]6.4}{[1][24][]24}}
\newlabel{ej:cv-digitos}{{6.12}{26}{Clasificador de dígitos}{ejemplo.6.12}{}}
\newlabel{ej:cv-digitos@cref}{{[ejemplo][12][]6.12}{[1][26][]26}}
\newlabel{eje:6-limites}{{6.1}{26}{}{ejercicio.6.1}{}}
\newlabel{eje:6-limites@cref}{{[ejercicio][1][6]6.1}{[1][26][]26}}
\newlabel{eje:6-ridge}{{6.2}{26}{}{ejercicio.6.2}{}}
\newlabel{eje:6-ridge@cref}{{[ejercicio][2][6]6.2}{[1][26][]26}}
\newlabel{eje:6-cv}{{6.3}{27}{}{ejercicio.6.3}{}}
\newlabel{eje:6-cv@cref}{{[ejercicio][3][6]6.3}{[1][26][]27}}
\newlabel{eje:6-lasso}{{6.4}{27}{}{ejercicio.6.4}{}}
\newlabel{eje:6-lasso@cref}{{[ejercicio][4][6]6.4}{[1][27][]27}}
\newlabel{eje:6-map}{{6.5}{27}{}{ejercicio.6.5}{}}
\newlabel{eje:6-map@cref}{{[ejercicio][5][6]6.5}{[1][27][]27}}
\newlabel{eje:6-ruido}{{6.6}{27}{}{ejercicio.6.6}{}}
\newlabel{eje:6-ruido@cref}{{[ejercicio][6][6]6.6}{[1][27][]27}}
\newlabel{ap:soluciones}{{A}{28}{Soluciones de los ejercicios}{appendix.A}{}}
\newlabel{ap:soluciones@cref}{{[appendix][1][2147483647]A}{[1][28][]28}}
\newlabel{sol:eje:1-tipos}{{A.1}{28}{Tema 1}{section*.3}{}}
\newlabel{sol:eje:1-tipos@cref}{{[subappendix][1][2147483647,1]A.1}{[1][28][]28}}
\newlabel{sol:eje:1-particion}{{A.1}{28}{Tema 1}{section*.4}{}}
\newlabel{sol:eje:1-particion@cref}{{[subappendix][1][2147483647,1]A.1}{[1][28][]28}}
\newlabel{sol:eje:1-bayes}{{A.1}{29}{Tema 1}{section*.5}{}}
\newlabel{sol:eje:1-bayes@cref}{{[subappendix][1][2147483647,1]A.1}{[1][29][]29}}
\newlabel{sol:eje:1-constante}{{A.1}{29}{Tema 1}{section*.6}{}}
\newlabel{sol:eje:1-constante@cref}{{[subappendix][1][2147483647,1]A.1}{[1][29][]29}}
\newlabel{sol:eje:1-logloss}{{A.1}{29}{Tema 1}{section*.7}{}}
\newlabel{sol:eje:1-logloss@cref}{{[subappendix][1][2147483647,1]A.1}{[1][29][]29}}
\newlabel{sol:eje:2-rango}{{A.2}{30}{Tema 2}{section*.8}{}}
\newlabel{sol:eje:2-rango@cref}{{[subappendix][2][2147483647,1]A.2}{[1][30][]30}}
\newlabel{sol:eje:2-recta}{{A.2}{31}{Tema 2}{section*.9}{}}
\newlabel{sol:eje:2-recta@cref}{{[subappendix][2][2147483647,1]A.2}{[1][30][]31}}
\newlabel{sol:eje:2-strang1}{{A.2}{31}{Tema 2}{section*.10}{}}
\newlabel{sol:eje:2-strang1@cref}{{[subappendix][2][2147483647,1]A.2}{[1][31][]31}}
\newlabel{sol:eje:2-strang2}{{A.2}{31}{Tema 2}{section*.11}{}}
\newlabel{sol:eje:2-strang2@cref}{{[subappendix][2][2147483647,1]A.2}{[1][31][]31}}
\newlabel{sol:eje:2-laplace}{{A.2}{32}{Tema 2}{section*.12}{}}
\newlabel{sol:eje:2-laplace@cref}{{[subappendix][2][2147483647,1]A.2}{[1][32][]32}}
\newlabel{sol:eje:2-traza}{{A.2}{32}{Tema 2}{section*.13}{}}
\newlabel{sol:eje:2-traza@cref}{{[subappendix][2][2147483647,1]A.2}{[1][32][]32}}
\newlabel{sol:eje:3-lineal}{{A.3}{33}{Tema 3}{section*.14}{}}
\newlabel{sol:eje:3-lineal@cref}{{[subappendix][3][2147483647,1]A.3}{[1][32][]33}}
\newlabel{sol:eje:3-valores}{{A.3}{33}{Tema 3}{section*.15}{}}
\newlabel{sol:eje:3-valores@cref}{{[subappendix][3][2147483647,1]A.3}{[1][33][]33}}
\newlabel{sol:eje:3-gradiente}{{A.3}{33}{Tema 3}{section*.16}{}}
\newlabel{sol:eje:3-gradiente@cref}{{[subappendix][3][2147483647,1]A.3}{[1][33][]33}}
\newlabel{sol:eje:3-frontera}{{A.3}{34}{Tema 3}{section*.17}{}}
\newlabel{sol:eje:3-frontera@cref}{{[subappendix][3][2147483647,1]A.3}{[1][34][]34}}
\newlabel{sol:eje:3-newton}{{A.3}{34}{Tema 3}{section*.18}{}}
\newlabel{sol:eje:3-newton@cref}{{[subappendix][3][2147483647,1]A.3}{[1][34][]34}}
\newlabel{sol:eje:3-separable}{{A.3}{35}{Tema 3}{section*.19}{}}
\newlabel{sol:eje:3-separable@cref}{{[subappendix][3][2147483647,1]A.3}{[1][35][]35}}
\newlabel{sol:eje:4-llm}{{A.4}{35}{Tema 4}{section*.20}{}}
\newlabel{sol:eje:4-llm@cref}{{[subappendix][4][2147483647,1]A.4}{[1][35][]35}}
\newlabel{sol:eje:4-logits}{{A.4}{36}{Tema 4}{section*.21}{}}
\newlabel{sol:eje:4-logits@cref}{{[subappendix][4][2147483647,1]A.4}{[1][36][]36}}
\newlabel{sol:eje:4-k2}{{A.4}{36}{Tema 4}{section*.22}{}}
\newlabel{sol:eje:4-k2@cref}{{[subappendix][4][2147483647,1]A.4}{[1][36][]36}}
\newlabel{sol:eje:4-llama}{{A.4}{36}{Tema 4}{section*.23}{}}
\newlabel{sol:eje:4-llama@cref}{{[subappendix][4][2147483647,1]A.4}{[1][36][]36}}
\newlabel{sol:eje:4-gradiente}{{A.4}{37}{Tema 4}{section*.24}{}}
\newlabel{sol:eje:4-gradiente@cref}{{[subappendix][4][2147483647,1]A.4}{[1][37][]37}}
\newlabel{sol:eje:4-hessiano}{{A.4}{37}{Tema 4}{section*.25}{}}
\newlabel{sol:eje:4-hessiano@cref}{{[subappendix][4][2147483647,1]A.4}{[1][37][]37}}
\newlabel{sol:eje:5-vf}{{A.5}{38}{Tema 5}{section*.26}{}}
\newlabel{sol:eje:5-vf@cref}{{[subappendix][5][2147483647,1]A.5}{[1][38][]38}}
\newlabel{sol:eje:5-diagnostico}{{A.5}{38}{Tema 5}{section*.27}{}}
\newlabel{sol:eje:5-diagnostico@cref}{{[subappendix][5][2147483647,1]A.5}{[1][38][]38}}
\newlabel{sol:eje:5-mano}{{A.5}{38}{Tema 5}{section*.28}{}}
\newlabel{sol:eje:5-mano@cref}{{[subappendix][5][2147483647,1]A.5}{[1][38][]38}}
\newlabel{sol:eje:5-interpolar}{{A.5}{39}{Tema 5}{section*.29}{}}
\newlabel{sol:eje:5-interpolar@cref}{{[subappendix][5][2147483647,1]A.5}{[1][39][]39}}
\newlabel{sol:eje:5-encoger}{{A.5}{39}{Tema 5}{section*.30}{}}
\newlabel{sol:eje:5-encoger@cref}{{[subappendix][5][2147483647,1]A.5}{[1][39][]39}}
\newlabel{sol:eje:5-varianza}{{A.5}{40}{Tema 5}{section*.31}{}}
\newlabel{sol:eje:5-varianza@cref}{{[subappendix][5][2147483647,1]A.5}{[1][40][]40}}
\newlabel{sol:eje:6-limites}{{A.6}{40}{Tema 6}{section*.32}{}}
\newlabel{sol:eje:6-limites@cref}{{[subappendix][6][2147483647,1]A.6}{[1][40][]40}}
\newlabel{sol:eje:6-ridge}{{A.6}{40}{Tema 6}{section*.33}{}}
\newlabel{sol:eje:6-ridge@cref}{{[subappendix][6][2147483647,1]A.6}{[1][40][]40}}
\newlabel{sol:eje:6-cv}{{A.6}{41}{Tema 6}{section*.34}{}}
\newlabel{sol:eje:6-cv@cref}{{[subappendix][6][2147483647,1]A.6}{[1][41][]41}}
\newlabel{sol:eje:6-lasso}{{A.6}{41}{Tema 6}{section*.35}{}}
\newlabel{sol:eje:6-lasso@cref}{{[subappendix][6][2147483647,1]A.6}{[1][41][]41}}
\newlabel{sol:eje:6-map}{{A.6}{41}{Tema 6}{section*.36}{}}
\newlabel{sol:eje:6-map@cref}{{[subappendix][6][2147483647,1]A.6}{[1][41][]41}}
\newlabel{sol:eje:6-ruido}{{A.6}{42}{Tema 6}{section*.37}{}}
\newlabel{sol:eje:6-ruido@cref}{{[subappendix][6][2147483647,1]A.6}{[1][42][]42}}
\makeatother

\begin{document}
\frontmatter
\begin{titlepage}
\centering
\vspace*{2.5cm}
{\large\scshape Anatomía de un LLM}\par
\vspace{0.4cm}
{\normalsize Fundamentos matemáticos de los modelos de lenguaje}\par
\vspace{2.2cm}
\rule{\textwidth}{0.8pt}\par\vspace{0.5cm}
{\small\scshape Fase II \textperiodcentered\ Aprendizaje automático y redes neuronales}\par\vspace{0.3cm}
{\Huge\bfseries Módulo 8\par}\vspace{0.3cm}
{\LARGE Aprendizaje supervisado\par}
\vspace{0.5cm}\rule{\textwidth}{0.4pt}\par
\vspace{1.6cm}
{\large Resumen del módulo}\par
\vfill
{\small Edición de octubre de 2026}\par
\end{titlepage}

\chapter*{Sobre este resumen}
\addcontentsline{toc}{chapter}{Sobre este resumen}
\markboth{Sobre este resumen}{}

Este documento reúne lo indispensable de los apuntes del Módulo 8: el texto expositivo, todas las definiciones y todos los resultados (teoremas, proposiciones, lemas y corolarios), las tablas y las observaciones de contenido. Se han suprimido las demostraciones, los ejemplos, las lecturas, las figuras, los ejercicios y sus soluciones. La numeración de temas, secciones, enunciados y ecuaciones es la misma que en los apuntes, de modo que cualquier referencia puede buscarse allí; las que remiten a elementos suprimidos (un ejemplo, una figura, una ecuación de una demostración) conservan su número en los apuntes.

\section*{Notación}
\begin{center}
\begin{tabular}{@{}ll@{}}
\toprule
Símbolo & Significado \\
\midrule
$(x_i,y_i)$, $i=1,\dots,n$ & ejemplos de entrenamiento; $x_i\in\R^{d}$ \\
$\D$ & distribución conjunta de $(x,y)$ \\
$X\in\R^{n\times(d+1)}$, $y\in\R^{n}$ & matriz de diseño y vector de respuestas \\
$\theta$, $w$, $b$ & parámetros, pesos y término independiente \\
$\ell(\hat y,y)$, $R(f)$, $\Remp(f)$ & pérdida, riesgo y riesgo empírico \\
$\sigma(z)=1/(1+e^{-z})$ & función sigmoide (logística) \\
$z\in\R^{K}$, $p=\softmax(z)$ & logits y probabilidades de las $K$ clases \\
$e_k$ & vector canónico (\emph{one-hot}) de la clase $k$ \\
$\ind[A]$ & indicador: vale 1 si $A$ es cierto y 0 si no \\
$\lambda$ & coeficiente de regularización \\
$\N(\mu,\sigma^{2})$ & distribución normal \\
$\ln$ & logaritmo natural; las pérdidas se miden en nats \\
\bottomrule
\end{tabular}
\end{center}


\tableofcontents

\mainmatter
% =====================================================================
\chapter{El problema del aprendizaje supervisado}
\label{tema:1}
% =====================================================================

Este tema fija el marco común a todo el módulo: los datos se suponen extraídos de una distribución desconocida, un modelo es una función elegida dentro de una clase de hipótesis y su calidad se mide con el riesgo, la pérdida esperada sobre datos nuevos. La presentación sigue la introducción de la parte~I de los apuntes de CS229~\cite{cs229}, la teoría de la decisión estadística de Hastie, Tibshirani y Friedman~\cite[§2.4]{esl}, la formulación del capítulo~2 de las notas de MIT~6.390~\cite{mit6390} y el análisis de las pérdidas de Bishop~\cite[§1.5.5]{bishop}.

\section{Datos, hipótesis y pérdida}

\begin{definicion}[Problema supervisado]\label{def:supervisado}
Sean $\X$ un \emph{espacio de entradas} y $\Y$ un \emph{espacio de salidas}. Un \emph{problema de aprendizaje supervisado} consta de una distribución de probabilidad $\D$ sobre $\X\times\Y$, desconocida, y de un \emph{conjunto de entrenamiento} $S=\{(x_1,y_1),\dots,(x_n,y_n)\}$ formado por $n$ pares independientes con distribución $\D$. El problema se llama de \emph{regresión} si $\Y=\R$ y de \emph{clasificación} si $\Y$ es finito; en este caso sus elementos se llaman \emph{clases}.
\end{definicion}

En casi todo el módulo $\X=\R^{d}$: cada entrada es un vector de $d$ \emph{características}. Un modelo de lenguaje preentrenado resuelve un problema de clasificación en el que $x$ es el contexto, $y$ es el token siguiente y $\Y=\V$ es el vocabulario, con $K=\abs{\V}$ clases.

\begin{definicion}[Hipótesis y clase de hipótesis]
Una \emph{hipótesis} es una función $f\colon\X\to\widehat{\Y}$, donde $\widehat{\Y}$ es el espacio de predicciones ($\widehat{\Y}=\R$ en regresión; $\widehat{\Y}=\Y$ o el símplex de distribuciones sobre $\Y$ en clasificación). Una \emph{clase de hipótesis} es un conjunto $\Hc=\{f_\theta:\theta\in\Theta\}$ de hipótesis indexadas por un vector de \emph{parámetros} $\theta$.
\end{definicion}

\begin{definicion}[Pérdida, riesgo y riesgo empírico]\label{def:riesgo}
Una \emph{función de pérdida} es una aplicación $\ell\colon\widehat{\Y}\times\Y\to[0,\infty)$; $\ell(\hat y,y)$ mide el coste de predecir $\hat y$ cuando la respuesta es $y$. Las tres pérdidas de este módulo son:
\begin{enumerate}
  \item la \emph{cuadrática}, $\ell(\hat y,y)=(\hat y-y)^{2}$, para regresión;
  \item la \emph{0-1}, $\ell(\hat y,y)=\ind[\hat y\ne y]$, para clasificación con predicción de una clase;
  \item la \emph{logarítmica}, $\ell(q,y)=-\ln q_y$, para clasificación con predicción de una distribución $q$ sobre $\Y$.
\end{enumerate}
El \emph{riesgo} o \emph{error de generalización} de $f$ es $R(f)=\E_{(x,y)\sim\D}\bigl[\ell(f(x),y)\bigr]$, y su \emph{riesgo empírico} sobre $S$ es
\[
\Remp(f)=\frac{1}{n}\sum_{i=1}^{n}\ell\bigl(f(x_i),y_i\bigr).
\]
El principio de \emph{minimización del riesgo empírico} elige $\hat f=\argmin_{f\in\Hc}\Remp(f)$.
\end{definicion}

El objetivo es un riesgo pequeño, pero $\D$ es desconocida y solo $\Remp$ puede calcularse. Toda la teoría de la generalización (temas~5 y~6) estudia la diferencia entre ambos.

\section{El predictor óptimo}

Antes de restringirse a una clase $\Hc$, conviene saber cuál es la mejor predicción posible si $\D$ fuese conocida~\cite[§2.4]{esl}.

\begin{teorema}[Función de regresión]\label{teo:regresion}
Sea $\Y=\R$ con $\E[y^{2}]<\infty$, y sea $f^{*}(x)=\E[y\mid x]$. Para toda $f$ con $\E[f(x)^{2}]<\infty$, el riesgo cuadrático cumple
\[
R(f)=\E\bigl[(y-f(x))^{2}\bigr]=\E\bigl[(y-f^{*}(x))^{2}\bigr]+\E\bigl[(f^{*}(x)-f(x))^{2}\bigr].
\]
En consecuencia, $f^{*}$ minimiza el riesgo cuadrático, y el riesgo mínimo es $\E\bigl[\Var(y\mid x)\bigr]$.
\end{teorema}

\begin{proposicion}[Clasificador de Bayes]\label{prop:bayes}
Sea $\Y$ finito y $\ell$ la pérdida 0-1. El clasificador $f^{*}(x)=\argmax_{k\in\Y}P(y=k\mid x)$ minimiza el riesgo, y su riesgo es $R(f^{*})=1-\E\bigl[\max_k P(y=k\mid x)\bigr]$.
\end{proposicion}

\begin{proposicion}[La pérdida logarítmica premia la probabilidad verdadera]\label{prop:logloss}
Sea $\Y$ finito, $p(\cdot\mid x)$ la distribución condicional verdadera y $q$ cualquier distribución sobre $\Y$. Entonces
\[
\E\bigl[-\ln q_y\mid x\bigr]=H\bigl(p(\cdot\mid x)\bigr)+D_{\mathrm{KL}}\bigl(p(\cdot\mid x)\,\Vert\,q\bigr),
\]
de modo que el riesgo logarítmico se minimiza, para cada $x$, con $q=p(\cdot\mid x)$.
\end{proposicion}

\addtocounter{definicion}{1}%

\addtocounter{definicion}{1}%

\begin{table}[h]
\centering\small
\caption{Distribución conjunta $P(x,y)$ del \cref{ej:bayes}.}\label{tab:conjunta}
\begin{tabular}{@{}lccc|c@{}}
\toprule
 & $y=0$ & $y=1$ & $y=2$ & $P(x)$ \\
\midrule
$x=0$ & $0{,}2$ & $0{,}1$ & $0{,}1$ & $0{,}4$ \\
$x=1$ & $0{,}1$ & $0{,}2$ & $0{,}3$ & $0{,}6$ \\
\bottomrule
\end{tabular}
\end{table}

\section{Riesgo empírico y optimismo}

\begin{proposicion}\label{prop:insesgado}
Sea $f$ una hipótesis fija, elegida sin mirar $S$, con $\Var(\ell(f(x),y))=v<\infty$. Entonces $\E_S[\Remp(f)]=R(f)$ y $\Var_S(\Remp(f))=v/n$.
\end{proposicion}

\begin{proposicion}[Optimismo del error de entrenamiento]\label{prop:optimismo}
Sea $\Hc$ una clase de hipótesis y supóngase que $\hat f\in\argmin_{f\in\Hc}\Remp(f)$ existe. Entonces
\[
\E_S\bigl[\Remp(\hat f)\bigr]\le\inf_{f\in\Hc}R(f)\le\E_S\bigl[R(\hat f)\bigr].
\]
\end{proposicion}

La \cref{prop:optimismo} explica por qué el error de entrenamiento no sirve para evaluar un modelo: en media, subestima el riesgo del modelo ajustado. Hace falta un conjunto de datos que no haya intervenido en el ajuste.

\section{Entrenamiento, validación y prueba}

\begin{definicion}[Partición de los datos]
Los datos etiquetados disponibles se dividen al azar en tres subconjuntos disjuntos: el de \emph{entrenamiento}, con el que se ajustan los parámetros $\theta$; el de \emph{validación}, con el que se eligen los \emph{hiperparámetros} (grado de un polinomio, coeficiente de regularización, número de épocas), y el de \emph{prueba}, con el que se estima una sola vez el riesgo del modelo final~\cite[§2.8]{mit6390}.
\end{definicion}

Si se elige un hiperparámetro mirando el conjunto de prueba, ese conjunto pasa a formar parte del ajuste y, por la \cref{prop:optimismo}, la estimación final resulta optimista. La \cref{fig:particion} resume el protocolo; la validación cruzada (\cref{tema:6}) lo refina cuando hay pocos datos.

\addtocounter{figure}{1}%

\addtocounter{definicion}{1}%


% =====================================================================
\chapter{Regresión lineal y mínimos cuadrados}
\label{tema:2}
% =====================================================================

Este tema estudia el modelo más sencillo de aprendizaje supervisado, la regresión lineal ajustada por mínimos cuadrados. Se deducen las ecuaciones normales, se interpretan como una proyección ortogonal y se justifica la pérdida cuadrática como máxima verosimilitud con ruido gaussiano. La exposición sigue el capítulo~1 de los apuntes de CS229~\cite{cs229}, las secciones~4.2 y~4.3 de Strang~\cite{strang} y el curso 18.06SC del MIT~\cite{mit1806}, la sección~3.1 de Bishop~\cite{bishop} y la sección~3.1 de Zhang \emph{et al.}~\cite{d2l}.

\section{El modelo lineal y la matriz de diseño}

\begin{definicion}[Regresión lineal]
Un \emph{modelo de regresión lineal} predice $\hat y=w^{\top}x+b$, con \emph{pesos} $w\in\R^{d}$ y \emph{término independiente} $b\in\R$. Añadiendo a cada entrada una componente constante, $\tilde x=(1,x)\in\R^{d+1}$, y escribiendo $\theta=(b,w)$, el modelo queda $\hat y=\theta^{\top}\tilde x$. En lo sucesivo se omite la tilde.
\end{definicion}

\begin{definicion}[Matriz de diseño y suma de cuadrados]
Dados $n$ ejemplos, la \emph{matriz de diseño} $X\in\R^{n\times(d+1)}$ tiene por fila $i$ el vector $x_i^{\top}$ (con su primera componente igual a~1), y $y=(y_1,\dots,y_n)^{\top}$. El vector de predicciones es $\hat y=X\theta$, y la \emph{suma de cuadrados de los residuos} es
\[
J(\theta)=\norm{y-X\theta}^{2}=\sum_{i=1}^{n}\bigl(y_i-\theta^{\top}x_i\bigr)^{2}.
\]
El estimador de \emph{mínimos cuadrados} es cualquier $\hat\theta$ que minimice $J$; se escribe $\SSE=J(\hat\theta)$.
\end{definicion}

\begin{observacion}[Funciones base]
El modelo es lineal en $\theta$, no necesariamente en $x$. Si se sustituye $x$ por un vector de \emph{funciones base} $\phi(x)$, por ejemplo $\phi(x)=(1,x,x^{2},\dots,x^{p})$, todo lo que sigue se aplica sin cambios con $X_{ij}=\phi_j(x_i)$~\cite[§3.1]{bishop}. Así se ajustan los polinomios del \cref{tema:5}.
\end{observacion}

\section{Las ecuaciones normales}

\begin{lema}[Gradientes de formas lineales y cuadráticas]\label{lem:gradientes}
Sean $a\in\R^{m}$ y $A\in\R^{m\times m}$. Entonces $\nabla_\theta(a^{\top}\theta)=a$ y $\nabla_\theta(\theta^{\top}A\theta)=(A+A^{\top})\theta$, que vale $2A\theta$ si $A$ es simétrica.
\end{lema}

\begin{teorema}[Ecuaciones normales]\label{teo:normales}
Un vector $\hat\theta$ minimiza $J(\theta)=\norm{y-X\theta}^{2}$ si y solo si cumple las \emph{ecuaciones normales}
\[
X^{\top}X\,\hat\theta=X^{\top}y.
\]
Además, para cualquier solución $\hat\theta$ y cualquier $\theta$, se verifica $J(\theta)=J(\hat\theta)+\norm{X(\theta-\hat\theta)}^{2}$.
\end{teorema}

\begin{lema}\label{lem:rango}
$X^{\top}X$ es simétrica y semidefinida positiva. Es definida positiva, y por tanto invertible, si y solo si las columnas de $X$ son linealmente independientes, es decir, si $\rg(X)=d+1$.
\end{lema}

\begin{corolario}[Unicidad]\label{cor:unicidad}
Si $\rg(X)=d+1$, el estimador de mínimos cuadrados es único y vale $\hat\theta=(X^{\top}X)^{-1}X^{\top}y$. Si $\rg(X)<d+1$, hay infinitos minimizadores, todos con las mismas predicciones $X\hat\theta$, y el de norma mínima es $X^{+}y$, donde $X^{+}$ es la pseudoinversa (módulo~3).
\end{corolario}

\addtocounter{definicion}{1}%

\begin{observacion}[Cálculo numérico]
En la práctica no se calcula $(X^{\top}X)^{-1}$: se resuelve el sistema con una factorización de Cholesky, o se aplica a $X$ una factorización QR o la descomposición en valores singulares, que son más estables porque el número de condición de $X^{\top}X$ es el cuadrado del de $X$. Con millones de parámetros, como en una red neuronal, no existe fórmula cerrada y se usa el descenso de gradiente (\cref{sec:lms}).
\end{observacion}

\section{Interpretación geométrica}

Las ecuaciones normales equivalen a $X^{\top}(y-X\hat\theta)=0$: el residuo es ortogonal a todas las columnas de $X$. Por tanto, $\hat y=X\hat\theta$ es el punto del espacio de columnas $\mathcal{C}(X)$ más próximo a $y$, su \emph{proyección ortogonal}~\cite[§4.2--4.3]{strang}.

\begin{proposicion}[Matriz de proyección]\label{prop:proyeccion}
Supóngase que $\rg(X)=d+1$ y sea $P=X(X^{\top}X)^{-1}X^{\top}$, de tamaño $n\times n$. Entonces:
\begin{enumerate}
  \item $\hat y=Py$ y $r=(I-P)y$;
  \item $P^{\top}=P$ y $P^{2}=P$;
  \item $\hat y^{\top}r=0$ y, en consecuencia, $\norm{y}^{2}=\norm{\hat y}^{2}+\norm{r}^{2}$;
  \item si $X$ contiene la columna de unos, $\sum_i r_i=0$, y la recta (o el hiperplano) ajustada pasa por el punto medio $(\bar x,\bar y)$.
\end{enumerate}
\end{proposicion}

\addtocounter{figure}{1}%

En el \cref{ej:recta}, $\norm{y}^{2}=46$, $\norm{\hat y}^{2}=45{,}8$ y $\norm{r}^{2}=0{,}2$, de acuerdo con la \cref{prop:proyeccion}(iii).

\section{Descenso de gradiente: el algoritmo LMS}
\label{sec:lms}

Cuando $n$ o $d$ son grandes, se minimiza $J$ iterativamente. Con la pérdida media $\frac{1}{2n}\norm{y-X\theta}^{2}$, el gradiente es $-\frac1n X^{\top}(y-X\theta)$, y el descenso de gradiente con tasa $\alpha>0$ actualiza
\[
\theta\leftarrow\theta+\frac{\alpha}{n}\sum_{i=1}^{n}\bigl(y_i-\theta^{\top}x_i\bigr)x_i.
\]
Si en cada paso se usa un solo ejemplo, se obtiene la \emph{regla LMS} (\emph{least mean squares}) o de Widrow-Hoff, $\theta\leftarrow\theta+\alpha\,(y_i-\theta^{\top}x_i)\,x_i$: la corrección es proporcional al error cometido en ese ejemplo~\cite[§1.1]{cs229}. Como $J$ es convexa (su hessiano $2X^{\top}X$ es semidefinido positivo, \cref{lem:rango}), el descenso de gradiente converge al mínimo global si la tasa es suficientemente pequeña (módulo~7).

\section{Interpretación probabilística}

La pérdida cuadrática no es una elección arbitraria: es la que resulta de suponer ruido gaussiano~\cite[§1.3]{cs229}, \cite[§3.1.1]{bishop}.

\begin{definicion}[Modelo lineal gaussiano]
Las respuestas cumplen $y_i=\theta^{\top}x_i+\varepsilon_i$, donde los errores $\varepsilon_1,\dots,\varepsilon_n$ son independientes con distribución $\N(0,\sigma^{2})$. Equivalentemente, $y_i\mid x_i\sim\N(\theta^{\top}x_i,\sigma^{2})$.
\end{definicion}

\begin{teorema}[Mínimos cuadrados y máxima verosimilitud]\label{teo:mv-gauss}
En el modelo lineal gaussiano, la log-verosimilitud negativa es
\[
-\ln L(\theta,\sigma^{2})=\frac{n}{2}\ln(2\pi\sigma^{2})+\frac{1}{2\sigma^{2}}\norm{y-X\theta}^{2}.
\]
Para cualquier $\sigma^{2}>0$, el estimador de máxima verosimilitud de $\theta$ es el de mínimos cuadrados, y el de $\sigma^{2}$ es $\hat\sigma^{2}=\SSE/n$.
\end{teorema}

En el \cref{ej:recta}, $\hat\sigma^{2}=0{,}2/4=0{,}05$. El razonamiento se generaliza: cada supuesto sobre la distribución de $y$ produce su pérdida de máxima verosimilitud (\cref{tab:mv}). Con ruido de Laplace, por ejemplo, la pérdida es la suma de valores absolutos (\cref{eje:2-laplace}).

\begin{table}[h]
\centering\small
\caption{Pérdida de máxima verosimilitud según la distribución supuesta para $y$.}\label{tab:mv}
\begin{tabular}{@{}llll@{}}
\toprule
Respuesta & Distribución & Pérdida & Modelo \\
\midrule
real & normal & cuadrática & regresión lineal (\cref{tema:2}) \\
binaria & Bernoulli & entropía cruzada binaria & regresión logística (\cref{tema:3}) \\
una de $K$ clases & categórica & entropía cruzada & regresión softmax (\cref{tema:4}) \\
token siguiente & categórica sobre $\V$ & entropía cruzada & capa de salida de un LLM \\
\bottomrule
\end{tabular}
\end{table}

\begin{observacion}[Física]
$\frac12\sum_i(y_i-\hat y_i)^{2}$ es la energía elástica de $n$ muelles verticales de constante unidad que unen cada dato con la recta. La recta de mínimos cuadrados es la posición de equilibrio, y las ecuaciones normales expresan que la fuerza total y el momento total de los muelles son nulos.
\end{observacion}


% =====================================================================
\chapter{Regresión logística}
\label{tema:3}
% =====================================================================

Para clasificar en dos clases no basta con una recta, cuyas predicciones recorren todo $\R$: hace falta transformar una puntuación lineal en una probabilidad. Este tema estudia la regresión logística, obtiene su pérdida como verosimilitud de Bernoulli, calcula su gradiente y su hessiano, demuestra su convexidad y presenta el método de Newton. Se sigue el capítulo~2 de los apuntes de CS229~\cite[§2.1 y §2.4]{cs229}, Bishop~\cite[§4.3.2--4.3.3]{bishop}, Hastie, Tibshirani y Friedman~\cite[§4.4.1]{esl} y las notas de MIT~6.390~\cite[§4.3--4.4]{mit6390}.

\section{La función sigmoide}

\begin{definicion}[Sigmoide y logit]\label{def:sigmoide}
La \emph{función sigmoide} o \emph{logística} es $\sigma\colon\R\to(0,1)$, $\sigma(z)=\dfrac{1}{1+e^{-z}}$. Su inversa es la función \emph{logit}, $\sigma^{-1}(p)=\ln\dfrac{p}{1-p}$, el logaritmo de la razón de probabilidades (\emph{log-odds}).
\end{definicion}

\begin{proposicion}[Propiedades de la sigmoide]\label{prop:sigmoide}
Para todo $z\in\R$:
\begin{enumerate}
  \item $\sigma(-z)=1-\sigma(z)$; en particular, $\sigma(0)=1/2$;
  \item $\sigma'(z)=\sigma(z)\bigl(1-\sigma(z)\bigr)\in(0,1/4]$, con máximo en $z=0$;
  \item $\sigma$ es estrictamente creciente, con $\lim_{z\to-\infty}\sigma(z)=0$ y $\lim_{z\to\infty}\sigma(z)=1$;
  \item $\ln\sigma(z)=-\ln(1+e^{-z})$ y $\ln(1-\sigma(z))=-\ln(1+e^{z})$.
\end{enumerate}
\end{proposicion}

\begin{observacion}[Física]
La ocupación media de un nivel de energía $E$ en la estadística de Fermi-Dirac es $1/\bigl(e^{(E-\mu)/(k_BT)}+1\bigr)=\sigma\bigl((\mu-E)/(k_BT)\bigr)$. La sigmoide es, por tanto, la distribución de Fermi-Dirac, del mismo modo que la softmax es la de Boltzmann (módulo~1).
\end{observacion}

\section{El modelo y la frontera de decisión}

\begin{definicion}[Regresión logística]\label{def:logistica}
Sea $\Y=\{0,1\}$. El modelo de \emph{regresión logística} con parámetros $\theta\in\R^{d+1}$ (incluido el término independiente) es
\[
P(y=1\mid x;\theta)=\sigma(\theta^{\top}x),\qquad P(y=0\mid x;\theta)=1-\sigma(\theta^{\top}x).
\]
El número $z=\theta^{\top}x$ es el \emph{logit} del ejemplo. El \emph{clasificador asociado} predice la clase~1 si $\sigma(z)\ge 1/2$.
\end{definicion}

\begin{proposicion}\label{prop:frontera}
En el modelo logístico, el logaritmo de la razón de probabilidades es lineal en $x$, $\ln\frac{P(y=1\mid x)}{P(y=0\mid x)}=\theta^{\top}x$, y la frontera de decisión $\{x:\sigma(\theta^{\top}x)=1/2\}$ es el hiperplano $\theta^{\top}x=0$.
\end{proposicion}

\section{Verosimilitud y entropía cruzada binaria}

\begin{proposicion}[Pérdida de la regresión logística]\label{prop:bce}
Sean $(x_i,y_i)$, $i=1,\dots,n$, ejemplos independientes con $y_i\mid x_i\sim\Ber(\sigma(\theta^{\top}x_i))$, y sea $p_i=\sigma(\theta^{\top}x_i)$. La log-verosimilitud negativa media es
\[
J(\theta)=\frac{1}{n}\sum_{i=1}^{n}\BCE(p_i,y_i),\qquad \BCE(p,y)=-\bigl[y\ln p+(1-y)\ln(1-p)\bigr],
\]
la \emph{entropía cruzada binaria}.
\end{proposicion}

\addtocounter{definicion}{1}%

\section{Gradiente, hessiano y convexidad}

\begin{teorema}[Gradiente y hessiano]\label{teo:grad-logistica}
Sea $X$ la matriz de diseño, $p=(p_1,\dots,p_n)$ con $p_i=\sigma(\theta^{\top}x_i)$ y $S=\diag\bigl(p_i(1-p_i)\bigr)$. Entonces
\[
\frac{\partial\,\BCE(p_i,y_i)}{\partial z_i}=p_i-y_i,\qquad
\nabla J(\theta)=\frac{1}{n}X^{\top}(p-y),\qquad
\nabla^{2}J(\theta)=\frac{1}{n}X^{\top}SX.
\]
El hessiano es semidefinido positivo, luego $J$ es convexa. Si $\rg(X)=d+1$, es definido positivo y $J$ es estrictamente convexa.
\end{teorema}

El gradiente tiene la forma «error por entrada», igual que en la regla LMS del \cref{tema:2}: si el modelo ya acierta con seguridad ($p_i\approx y_i$), el ejemplo apenas contribuye. Esta coincidencia no es casual: ambos son modelos lineales generalizados con función de enlace canónica~\cite[cap.~3]{cs229}, \cite[§4.3.6]{bishop}.

\begin{corolario}[Calibración en media]\label{cor:calibracion}
Si $X$ contiene la columna de unos y $\hat\theta$ minimiza $J$, entonces $\sum_i p_i=\sum_i y_i$: la probabilidad media predicha coincide con la frecuencia observada de la clase~1.
\end{corolario}

A diferencia de las ecuaciones normales, la condición $X^{\top}(p-y)=0$ no es lineal en $\theta$, porque $p$ depende de $\theta$ a través de la sigmoide, y no tiene solución en forma cerrada. Puede incluso no tener solución.

\begin{proposicion}[Datos separables]\label{prop:separable}
Supóngase que existe $\theta_0$ con $\theta_0^{\top}x_i>0$ si $y_i=1$ y $\theta_0^{\top}x_i<0$ si $y_i=0$. Entonces $\inf_\theta J(\theta)=0$, pero el ínfimo no se alcanza: el estimador de máxima verosimilitud no existe.
\end{proposicion}

Con datos separables, el descenso de gradiente hace crecer $\norm{\theta}$ sin límite y las probabilidades predichas se acercan a 0 o~1: una forma de sobreajuste que Bishop señala en~\cite[§4.3.2]{bishop}. La regularización del \cref{tema:6} lo corrige.

\section{El método de Newton}

Como $J$ es convexa y dos veces diferenciable, puede minimizarse con el \emph{método de Newton}, que en cada paso minimiza la aproximación de Taylor de segundo orden (módulo~4):
\[
\theta_{t+1}=\theta_t-\bigl(\nabla^{2}J(\theta_t)\bigr)^{-1}\nabla J(\theta_t)=\theta_t-\bigl(X^{\top}SX\bigr)^{-1}X^{\top}(p-y).
\]
Cerca del óptimo, el método converge cuadráticamente: el número de cifras correctas aproximadamente se duplica en cada iteración~\cite[§2.4]{cs229}. Cada paso exige resolver un sistema de tamaño $d+1$, con un coste del orden de $d^{3}$, razón por la que no se usa en redes con millones de parámetros. Cada paso equivale a resolver un problema de mínimos cuadrados con pesos $S$, de ahí el nombre de \emph{mínimos cuadrados iterativamente reponderados} (IRLS)~\cite[§4.3.3]{bishop}.


% =====================================================================
\chapter{Regresión softmax y la capa de salida de un LLM}
\label{tema:4}
% =====================================================================

La regresión softmax generaliza la logística a $K$ clases: cada ejemplo recibe un vector de $K$ logits, que la función softmax convierte en una distribución. Este tema define el modelo, demuestra que el gradiente de su entropía cruzada respecto a los logits es $p-e_y$, calcula su hessiano y muestra que la última capa de un gran modelo de lenguaje es exactamente este modelo. Se sigue la sección~2.3 de los apuntes de CS229~\cite{cs229}, la sección~4.1 de Zhang \emph{et al.}~\cite{d2l}, la sección~4.3.4 de Bishop~\cite{bishop} y, para la conexión con los LLM, los artículos de Llama~2 y Llama~3~\cite{llama2,llama3}.

\section{El modelo}

\begin{definicion}[Regresión softmax]\label{def:softmax-reg}
Sea $\Y=\{1,\dots,K\}$. El modelo de \emph{regresión softmax} (o regresión logística multinomial) con parámetros $W\in\R^{d\times K}$ y $b\in\R^{K}$ asigna a cada entrada $x\in\R^{d}$ el vector de logits $z=W^{\top}x+b\in\R^{K}$ y la distribución
\[
P(y=k\mid x)=p_k=\softmax(z)_k=\frac{e^{z_k}}{\sum_{j=1}^{K}e^{z_j}},\qquad k=1,\dots,K.
\]
\end{definicion}

Si $w_k$ es la columna $k$ de $W$, el logit $z_k=w_k^{\top}x+b_k$ es el producto escalar de la entrada con una \emph{plantilla} de la clase $k$ (módulo~2): el modelo asigna más probabilidad a la clase cuya plantilla se parece más a la entrada. En la práctica del módulo, las columnas de $W$ entrenadas sobre imágenes de $8\times 8$ píxeles se pueden dibujar y recuerdan a los dígitos.

\begin{proposicion}[Relación con la regresión logística]\label{prop:k2}
\begin{enumerate}
  \item Si se suma un mismo vector $c$ a todas las columnas de $W$ y una misma constante a todas las componentes de $b$, las probabilidades no cambian. El modelo tiene, por tanto, parámetros redundantes.
  \item Con $K=2$, $p_1=\sigma(z_1-z_2)$: la regresión logística es la regresión softmax con dos clases y logit $z_1-z_2$.
\end{enumerate}
\end{proposicion}

\section{La entropía cruzada}

\begin{proposicion}[Pérdida de la regresión softmax]\label{prop:ce}
Si $y\mid x\sim\Cat(p_1,\dots,p_K)$ con $p=\softmax(z)$, la log-verosimilitud negativa de un ejemplo de clase $y$ es
\[
\ell(z,y)=-\ln p_y=-z_y+\LSE(z),\qquad \LSE(z)=\ln\sum_{j=1}^{K}e^{z_j},
\]
que coincide con la entropía cruzada $H(e_y,p)$ entre la distribución degenerada en la clase observada y la predicha.
\end{proposicion}

\addtocounter{definicion}{1}%

\section{El gradiente}

\begin{teorema}[Gradiente de la entropía cruzada]\label{teo:grad-softmax}
Sea $p=\softmax(z)$ y $\ell(z,y)=-\ln p_y$. Entonces
\[
\nabla_z\ell=p-e_y,\qquad\text{y en particular}\quad\sum_{k=1}^{K}\frac{\partial\ell}{\partial z_k}=0.
\]
\end{teorema}

La suma nula refleja la invariancia por traslación: desplazar todos los logits en la misma cantidad no cambia la pérdida, luego la derivada en la dirección $(1,\dots,1)$ es cero. Subir el logit de la clase correcta equivale a bajar los demás.

\begin{corolario}[Gradiente respecto a los parámetros]\label{cor:grad-W}
Para un ejemplo, $\nabla_W\ell=x\,(p-e_y)^{\top}\in\R^{d\times K}$ y $\nabla_b\ell=p-e_y$. Para la pérdida media sobre $n$ ejemplos, con $X\in\R^{n\times d}$, $P\in\R^{n\times K}$ la matriz de probabilidades predichas e $Y\in\R^{n\times K}$ la de etiquetas \emph{one-hot},
\[
\nabla_W J=\frac{1}{n}X^{\top}(P-Y),\qquad \nabla_bJ=\frac1n(P-Y)^{\top}\mathbf{1}.
\]
\end{corolario}

En el \cref{ej:softmax}, $\nabla_z\ell=(-0{,}3348;\,0{,}2447;\,0{,}0900)$: el descenso de gradiente sube el primer logit y baja los otros dos, y las componentes suman cero.

\section{Hessiano y convexidad}

\begin{proposicion}[Hessiano respecto a los logits]\label{prop:hess-softmax}
El hessiano de $\ell(z,y)$ respecto a $z$ es $\nabla_z^{2}\ell=\diag(p)-pp^{\top}$, que no depende de $y$. Es la matriz de covarianza de un vector aleatorio \emph{one-hot} $E$ con $P(E=e_k)=p_k$; por tanto, es semidefinido positivo, y su núcleo contiene el vector $\mathbf{1}$.
\end{proposicion}

\begin{corolario}\label{cor:convexa-softmax}
La pérdida media de la regresión softmax es una función convexa de $(W,b)$.
\end{corolario}

La convexidad no es estricta, por la redundancia de la \cref{prop:k2}(i): hay una dirección en la que la pérdida es constante. Aun así, todo mínimo local es global, y el descenso de gradiente no queda atrapado en mínimos locales malos. Esta garantía se pierde en cuanto la entrada $x$ la produce una red profunda (módulo~9).

\section{La capa de salida de un LLM}

Un modelo de lenguaje de tipo Transformer procesa el contexto con una pila de bloques y obtiene, para la última posición, un vector oculto $h\in\R^{d_{\text{model}}}$. La capa final, llamada \emph{unembedding} o \emph{cabeza de lenguaje}, calcula los logits $z=W_U^{\top}h$ con $W_U\in\R^{d_{\text{model}}\times\abs{\V}}$, y la softmax los convierte en la distribución del token siguiente (\cref{fig:cabeza}). Es una regresión softmax con $K=\abs{\V}$ clases cuya entrada $h$ no es fija, sino que la aprende la red. El preentrenamiento minimiza la entropía cruzada de la \cref{prop:ce} sobre cada posición de cada texto, y el gradiente que recibe la red a través de los logits es el $p-e_y$ del \cref{teo:grad-softmax}.

\addtocounter{figure}{1}%

\begin{proposicion}[Pérdida inicial]\label{prop:inicial}
Si $W=0$ y $b=0$, o en general si todos los logits son iguales, la pérdida de cualquier ejemplo es $\ln K$.
\end{proposicion}

El valor $\ln K$ es la referencia de un modelo que no ha aprendido nada: $\ln 10=2{,}303$ nats para los dígitos y $\ln 128\,256=11{,}76$ nats para un vocabulario como el de Llama~3. El \cref{tab:cabezas} compara el tamaño de la capa en ambos casos con datos publicados por sus autores~\cite{llama2,llama3}. Anthropic, Google y xAI no publican $d_{\text{model}}$ ni el tamaño del vocabulario de Claude, Gemini o Grok; lo que sí es público es que se trata de Transformers que terminan en esta capa.

\begin{table}[h]
\centering\small
\caption{Parámetros de la capa de salida (sin término independiente en los LLM).}\label{tab:cabezas}
\begin{tabular}{@{}lrrr@{}}
\toprule
Modelo & Entrada $d$ & Clases $K$ & Parámetros \\
\midrule
Clasificador de dígitos de la práctica & $64$ & $10$ & $64\cdot 10+10=650$ \\
Llama~2 7B~\cite{llama2} & $4096$ & $32\,000$ & $131\,072\,000$ \\
Llama~3 8B~\cite{llama3} & $4096$ & $128\,256$ & $525\,336\,576$ \\
\bottomrule
\end{tabular}
\end{table}

\begin{observacion}[Modelos lineales generalizados]
La regresión lineal, la logística y la softmax son casos de un mismo esquema, el de los \emph{modelos lineales generalizados}: la distribución de $y$ pertenece a la familia exponencial (normal, Bernoulli, categórica) y su parámetro natural es una función lineal de $x$. Con esa elección (el enlace canónico), el gradiente de la log-verosimilitud negativa es siempre «predicción menos observación» por la entrada~\cite[cap.~3]{cs229}, \cite[§4.3.6]{bishop}. Las fórmulas $(\hat y-y)x$, $(p-y)x$ y $x(p-e_y)^{\top}$ de los temas~2 a~4 son tres instancias del mismo resultado.
\end{observacion}


% =====================================================================
\chapter{Generalización: sesgo y varianza}
\label{tema:5}
% =====================================================================

Un modelo que reproduce bien los datos de entrenamiento puede predecir mal datos nuevos. Este tema estudia ese fenómeno, el sobreajuste, y lo explica mediante la descomposición del error en sesgo, varianza y ruido. Se sigue el capítulo~8 de los apuntes de CS229~\cite[§8.1--8.2]{cs229}, las secciones~2.9, 7.2 y~7.3 de Hastie, Tibshirani y Friedman~\cite{esl}, las secciones~1.1 y~3.2 de Bishop~\cite{bishop} y la sección~3.6 de Zhang \emph{et al.}~\cite{d2l}.

\section{Error de entrenamiento y error de generalización}

El experimento clásico de Bishop~\cite[§1.1]{bishop} consiste en ajustar polinomios de grado creciente a pocos puntos ruidosos de una función suave. En la práctica del módulo se toman $n=10$ puntos $x_i$ uniformes en $[0,1]$, con $y_i=\sin(2\pi x_i)+\varepsilon_i$ y $\varepsilon_i\sim\N(0;\,0{,}2^{2})$, se ajusta por mínimos cuadrados un polinomio de grado $p$ (funciones base $1,x,\dots,x^{p}$, \cref{tema:2}) y se mide el error cuadrático medio (ECM) sobre los datos de entrenamiento y sobre 200 puntos nuevos. El \cref{tab:polinomios} recoge el resultado.

\begin{table}[h]
\centering\small
\caption{Ajuste polinómico de 10 puntos de $\sin(2\pi x)$ con ruido. La última fila usa la regresión ridge del \cref{tema:6}.}\label{tab:polinomios}
\begin{tabular}{@{}cccccl@{}}
\toprule
Grado $p$ & $\lambda$ & Parámetros & ECM entrenamiento & ECM validación & Diagnóstico \\
\midrule
1 & 0 & 2 & $0{,}1183$ & $0{,}2888$ & infraajuste \\
3 & 0 & 4 & $0{,}0024$ & $0{,}0419$ & equilibrio \\
9 & 0 & 10 & $0{,}0000$ & $34{,}17$ & sobreajuste \\
9 & $0{,}001$ & 10 & $0{,}0144$ & $0{,}0652$ & regularizado \\
\bottomrule
\end{tabular}
\end{table}

Con $p=9$ hay tantos parámetros como datos y el polinomio pasa exactamente por los diez puntos (\cref{prop:interpolacion}), pero oscila violentamente entre ellos: sus coeficientes alcanzan valores del orden de $5\cdot 10^{6}$.

\begin{definicion}[Infraajuste y sobreajuste]
Un modelo \emph{infraajusta} si su error de entrenamiento es alto: la clase de hipótesis no puede representar la relación entre $x$ e $y$. Un modelo \emph{sobreajusta} si su error de entrenamiento es pequeño y su error de validación es mucho mayor: el modelo ha reproducido el ruido particular de la muestra.
\end{definicion}

\begin{proposicion}[Clases encajadas]\label{prop:encajadas}
Si $\Hc_1\subseteq\Hc_2$, el mínimo del riesgo empírico sobre $\Hc_2$ es menor o igual que sobre $\Hc_1$. En particular, el error de entrenamiento de mínimos cuadrados no aumenta al subir el grado del polinomio.
\end{proposicion}

\begin{proposicion}[Interpolación]\label{prop:interpolacion}
Si $x_1,\dots,x_n$ son distintos, existe un único polinomio de grado a lo sumo $n-1$ que pasa por los $n$ puntos $(x_i,y_i)$; su error de entrenamiento es nulo.
\end{proposicion}

La \cref{prop:encajadas} explica por qué el error de entrenamiento no sirve para elegir el grado: siempre favorece al modelo más flexible. El \cref{tab:polinomios} muestra que el error de validación, en cambio, primero baja y después sube.

\section{La descomposición sesgo-varianza}

Para entender la forma en U del error de validación, se considera el conjunto de entrenamiento $S$ como aleatorio y se estudia cómo varía la predicción $\hat f_S(x)$ en un punto fijo $x$~\cite[§8.1.1]{cs229}, \cite[§7.3]{esl}.

\begin{teorema}[Descomposición sesgo-varianza]\label{teo:sesgo-varianza}
Sea $y=f(x)+\varepsilon$, donde $f$ es determinista y $\varepsilon$ tiene media nula y varianza $\sigma^{2}$. Sea $\hat f_S$ un modelo entrenado con un conjunto aleatorio $S$, independiente de $\varepsilon$, y sea $\bar f(x)=\E_S[\hat f_S(x)]$. Para cada $x$ fijo,
\[
\E_{S,\varepsilon}\bigl[(y-\hat f_S(x))^{2}\bigr]
=\underbrace{\sigma^{2}}_{\text{ruido}}
+\underbrace{\bigl(f(x)-\bar f(x)\bigr)^{2}}_{\text{sesgo}^{2}}
+\underbrace{\E_S\bigl[(\hat f_S(x)-\bar f(x))^{2}\bigr]}_{\text{varianza}}.
\]
\end{teorema}

\begin{definicion}[Sesgo y varianza de un modelo]
En las condiciones del \cref{teo:sesgo-varianza}, el \emph{sesgo} en $x$ es $\bar f(x)-f(x)$, el error sistemático del modelo promediado sobre los posibles conjuntos de entrenamiento, y la \emph{varianza} es $\Var_S(\hat f_S(x))$, la sensibilidad de la predicción al conjunto concreto. El \emph{ruido} $\sigma^{2}$ es el error irreducible: ningún modelo puede eliminarlo (\cref{teo:regresion}).
\end{definicion}

Un modelo demasiado simple tiene sesgo alto (no puede seguir a $f$) y varianza baja; uno demasiado flexible tiene sesgo bajo y varianza alta. El error total se minimiza con una complejidad intermedia, lo que se conoce como \emph{compromiso sesgo-varianza}~\cite[§2.9]{esl}.

\addtocounter{definicion}{1}%

\begin{table}[h]
\centering\small
\caption{Sesgo y varianza del estimador $a\bar x$ del \cref{ej:encoger}.}\label{tab:encoger}
\begin{tabular}{@{}lccc@{}}
\toprule
$a$ & Sesgo$^{2}$ & Varianza & ECM \\
\midrule
$1$ (media muestral) & $0$ & $1$ & $1{,}00$ \\
$0{,}8$ & $0{,}16$ & $0{,}64$ & $0{,}80$ \\
$0{,}5$ & $1$ & $0{,}25$ & $1{,}25$ \\
\bottomrule
\end{tabular}
\end{table}

\begin{proposicion}[Varianza de los mínimos cuadrados]\label{prop:var-mc}
En el modelo lineal gaussiano con diseño $X$ fijo de rango $m$, el estimador de mínimos cuadrados es insesgado, $\E[\hat\theta]=\theta$, y $\Cov(\hat\theta)=\sigma^{2}(X^{\top}X)^{-1}$.
\end{proposicion}

Cada parámetro adicional aporta varianza; el \cref{eje:5-varianza} cuantifica la aportación sobre los puntos de entrenamiento.

\section{Diagnóstico y remedios}

El \cref{tab:diagnostico} traduce la descomposición en reglas prácticas.

\begin{table}[h]
\centering\small
\caption{Diagnóstico a partir de los errores de entrenamiento y validación.}\label{tab:diagnostico}
\begin{tabular}{@{}>{\raggedright\arraybackslash}p{0.34\textwidth}>{\raggedright\arraybackslash}p{0.15\textwidth}>{\raggedright\arraybackslash}p{0.42\textwidth}@{}}
\toprule
Síntoma & Diagnóstico & Remedios \\
\midrule
Error alto en entrenamiento y en validación & sesgo alto & modelo más rico, más características, menos regularización, entrenar más \\
Error bajo en entrenamiento y alto en validación & varianza alta & más datos, más regularización, modelo más simple, parada temprana \\
\bottomrule
\end{tabular}
\end{table}

\begin{observacion}[Doble descenso]
En modelos muy sobreparametrizados, como las redes neuronales, se observa que, tras el pico de error que aparece cuando el número de parámetros se acerca al de datos, el error de validación vuelve a bajar al seguir aumentando el tamaño del modelo. El fenómeno, llamado \emph{doble descenso}, no contradice el \cref{teo:sesgo-varianza}, pero muestra que el número de parámetros no es siempre una buena medida de la complejidad efectiva~\cite[§8.2]{cs229}.
\end{observacion}

\addtocounter{definicion}{1}%


% =====================================================================
\chapter{Regularización y selección de modelos}
\label{tema:6}
% =====================================================================

Regularizar es añadir a la pérdida un término que penaliza los modelos complicados, para cambiar un poco de sesgo por una reducción mayor de la varianza. Este tema estudia la regresión ridge, su interpretación bayesiana, el lasso y el \emph{weight decay}, y la validación cruzada con la que se eligen los hiperparámetros. Se sigue el capítulo~9 de los apuntes de CS229~\cite[§9.1--9.4]{cs229}, las secciones~3.4 y~7.10 de Hastie, Tibshirani y Friedman~\cite{esl}, las secciones~2.6 a~2.8 de las notas de MIT~6.390~\cite{mit6390}, la sección~3.1.4 de Bishop~\cite{bishop} y la sección~3.7 de Zhang \emph{et al.}~\cite{d2l}. Los métodos proceden de los artículos de Hoerl y Kennard~\cite{hoerl} y de Tibshirani~\cite{lasso}.

\section{Regresión ridge}

\begin{definicion}[Regresión ridge]\label{def:ridge}
Sea $\lambda>0$. El estimador \emph{ridge} es el minimizador de
\[
J_\lambda(\theta)=\norm{y-X\theta}^{2}+\lambda\norm{\theta}^{2}.
\]
\end{definicion}

\begin{teorema}\label{teo:ridge}
Para todo $\lambda>0$ y toda matriz $X$, $X^{\top}X+\lambda I$ es definida positiva, y el estimador ridge existe, es único y vale
\[
\hat\theta_\lambda=(X^{\top}X+\lambda I)^{-1}X^{\top}y.
\]
\end{teorema}

A diferencia de los mínimos cuadrados, ridge no necesita que $X$ tenga rango completo: resuelve también los problemas con más parámetros que datos. Esta fue la motivación original de Hoerl y Kennard~\cite{hoerl}.

\begin{proposicion}[Contracción en la base de valores singulares]\label{prop:ridge-svd}
Sea $X=UDV^{\top}$ la descomposición en valores singulares reducida, con $u_j$ las columnas de $U$ y $d_j>0$ los valores singulares. Entonces
\[
X\hat\theta_\lambda=\sum_{j}\frac{d_j^{2}}{d_j^{2}+\lambda}\,u_j\,(u_j^{\top}y).
\]
\end{proposicion}

Los mínimos cuadrados ($\lambda=0$) proyectan $y$ sobre cada dirección $u_j$ sin modificarla (\cref{prop:proyeccion}); ridge multiplica cada componente por $d_j^{2}/(d_j^{2}+\lambda)<1$ y contrae más las direcciones de valor singular pequeño, que son las que el ruido perturba más~\cite[§3.4.1]{esl}. La traza de la matriz que lleva $y$ a $X\hat\theta_\lambda$, $\mathrm{df}(\lambda)=\sum_j d_j^{2}/(d_j^{2}+\lambda)$, se llama \emph{número efectivo de parámetros}: vale $m$ con $\lambda=0$ (\cref{eje:2-traza}) y tiende a 0 cuando $\lambda\to\infty$.

\addtocounter{definicion}{1}%

\begin{observacion}[El término independiente no se penaliza]
El término independiente solo desplaza las predicciones y no añade flexibilidad, por lo que no suele penalizarse. Si se centran las columnas de $X$ y el vector $y$, el término independiente óptimo es nulo en el problema centrado, la penalización recae solo sobre $w$ y, al final, $b=\bar y-\hat w^{\top}\bar x$~\cite[§2.6 y §2.7.2]{mit6390}, \cite[§3.4.1]{esl}. Además, ridge no es invariante frente a cambios de escala de las características, que por eso se estandarizan antes de ajustar.
\end{observacion}

\section{Interpretación bayesiana}

\begin{proposicion}[Ridge como estimación MAP]\label{prop:map}
Sea $y\mid\theta\sim\N(X\theta,\sigma^{2}I)$ y sea la distribución \emph{a priori} $\theta\sim\N(0,\tau^{2}I)$. El estimador de máximo \emph{a posteriori} (MAP) de $\theta$ es el estimador ridge con $\lambda=\sigma^{2}/\tau^{2}$.
\end{proposicion}

La demostración se propone en el \cref{eje:6-map}. Regularizar equivale a creer, antes de ver los datos, que los pesos son pequeños; cuanto más firme es esa creencia (menor $\tau$), mayor es $\lambda$~\cite[§9.4]{cs229}, \cite[§3.1.4]{bishop}. Con una distribución \emph{a priori} de Laplace, $p(\theta_j)\propto e^{-\abs{\theta_j}/s}$, el mismo razonamiento conduce a la penalización $\norm{\theta}_1$ de la sección siguiente.

\section{El lasso}

\begin{definicion}[Lasso]
El estimador \emph{lasso}~\cite{lasso} minimiza $\frac12\norm{y-X\theta}^{2}+\lambda\norm{\theta}_1$, con $\norm{\theta}_1=\sum_j\abs{\theta_j}$.
\end{definicion}

\begin{proposicion}[Umbral suave]\label{prop:umbral}
Si las columnas de $X$ son ortonormales ($X^{\top}X=I$) y $\hat\beta=X^{\top}y$ es el estimador de mínimos cuadrados, el lasso tiene solución única
\[
\hat\theta_j=\sign(\hat\beta_j)\,\bigl(\abs{\hat\beta_j}-\lambda\bigr)_{+},
\]
donde $(t)_{+}=\max(t,0)$. Con la penalización $\frac{\lambda}{2}\norm{\theta}^{2}$ en lugar de $\lambda\norm{\theta}_1$, la solución es $\hat\theta_j=\hat\beta_j/(1+\lambda)$.
\end{proposicion}

Ridge contrae todos los coeficientes en la misma proporción y nunca los anula; el lasso resta una cantidad fija y anula los que no superan el umbral, de modo que selecciona variables y produce modelos \emph{dispersos}~\cite[§3.4.2--3.4.3]{esl}. Las penalizaciones de tipo $\ell_1$ reaparecen en los autocodificadores dispersos de la interpretabilidad (módulo~23).

\section{\emph{Weight decay} y otras formas de regularización}

\begin{proposicion}[Weight decay]\label{prop:weight-decay}
Un paso de descenso de gradiente con tasa $\eta$ sobre $J(\theta)+\frac{\lambda}{2}\norm{\theta}^{2}$ es
\[
\theta\leftarrow(1-\eta\lambda)\,\theta-\eta\,\nabla J(\theta):
\]
antes del paso habitual, cada peso se multiplica por el factor $1-\eta\lambda<1$.
\end{proposicion}

Con $\eta=0{,}1$ y $\lambda=0{,}01$, el factor es $0{,}999$: los pesos decaen geométricamente salvo que el gradiente de los datos los sostenga~\cite[§3.7.1]{d2l}. La práctica del módulo usa precisamente la pérdida $J+\frac{\lambda}{2}\norm{W}^{2}$, con gradiente $\frac1nX^{\top}(P-Y)+\lambda W$. Con optimizadores adaptativos como Adam la equivalencia entre penalización $\ell_2$ y decaimiento deja de cumplirse, y AdamW aplica el decaimiento por separado~\cite{adamw} (módulo~7).

\begin{table}[h]
\centering\small
\caption{Técnicas de regularización y dónde aparecen en el curso.}\label{tab:regularizacion}
\begin{tabular}{@{}>{\raggedright\arraybackslash}p{0.27\textwidth}>{\raggedright\arraybackslash}p{0.36\textwidth}>{\raggedright\arraybackslash}p{0.29\textwidth}@{}}
\toprule
Técnica & Efecto & En el curso \\
\midrule
$\ell_2$, ridge, \emph{weight decay} & contrae todos los pesos & AdamW (módulo~7), casi todos los LLM \\
$\ell_1$, lasso & anula pesos; dispersión & autocodificadores dispersos (módulo~23) \\
Parada temprana & detiene el ajuste cuando sube la validación & ajuste fino (módulo~20) \\
\emph{Dropout} & apaga unidades al azar durante el entrenamiento & módulo~11 \\
Más datos & reduce directamente la varianza & estrategia principal de los LLM \\
\bottomrule
\end{tabular}
\end{table}

La parada temprana y el propio descenso de gradiente estocástico tienen un efecto de \emph{regularización implícita}: entre las muchas soluciones de un problema sobreparametrizado, el algoritmo tiende a encontrar unas concretas, a menudo de norma pequeña~\cite[§9.2]{cs229}.

\section{Selección de modelos y validación cruzada}

\begin{definicion}[Validación cruzada con $k$ particiones]\label{def:kfold}
Sea $S$ el conjunto de entrenamiento y $k\ge 2$. Se divide $S$ al azar en $k$ partes disjuntas $S_1,\dots,S_k$ de tamaños lo más parecidos posible. Para cada valor del hiperparámetro y cada $j=1,\dots,k$, se entrena el modelo con $S\setminus S_j$ y se mide su error en $S_j$. La \emph{estimación por validación cruzada} es la media de los $k$ errores. Se elige el hiperparámetro con menor estimación y, por último, se reentrena con todo $S$~\cite[§9.3]{cs229}, \cite[§7.10.1]{esl}.
\end{definicion}

Cada ejemplo se usa para validar exactamente una vez por cada valor del hiperparámetro, y la media de $k$ errores tiene menos varianza que el error en una única partición de validación. El precio es el coste: $k$ entrenamientos por cada valor. El caso $k=n$ se llama \emph{dejar uno fuera} (\emph{leave-one-out}); su estimación es casi insesgada, pero su coste es de $n$ entrenamientos.

\begin{observacion}[Uso correcto]
Todo paso que mire las etiquetas, como seleccionar características o elegir umbrales, debe repetirse dentro de cada partición, usando solo $S\setminus S_j$; si se hace antes con todo $S$, la estimación resulta optimista~\cite[§7.10.2]{esl}. Con series temporales no se baraja: se valida siempre con datos posteriores a los de entrenamiento.
\end{observacion}

\addtocounter{definicion}{1}%

\begin{table}[h]
\centering\small
\caption{Exactitud media por validación cruzada ($k=5$) del \cref{ej:cv-digitos}.}\label{tab:cv}
\begin{tabular}{@{}lcccc@{}}
\toprule
$\lambda$ & $0$ & $0{,}001$ & $0{,}01$ & $0{,}1$ \\
\midrule
Exactitud & $0{,}949$ & $0{,}950$ & $0{,}937$ & $0{,}719$ \\
\bottomrule
\end{tabular}
\end{table}

\addtocounter{definicion}{1}%



\backmatter
% =====================================================================
\begin{thebibliography}{99}
\addcontentsline{toc}{chapter}{Bibliografía}
\markboth{Bibliografía}{}
\small

\bibitem{bishop}
C.~M. Bishop.
\newblock \emph{Pattern Recognition and Machine Learning}.
\newblock Springer, 2006.
\newblock Caps.~1, 3 y~4.
\newblock \url{https://www.microsoft.com/en-us/research/people/cmbishop/prml-book/}

\bibitem{llama3}
A.~Grattafiori \emph{et al.} (Llama Team, AI @ Meta).
\newblock The Llama 3 herd of models.
\newblock arXiv:2407.21783, 2024.
\newblock Los valores $d_{\text{model}}=4096$ y $\abs{\V}=128\,256$ de Llama~3 8B figuran en el archivo de configuración publicado con los pesos del modelo.
\newblock \url{https://arxiv.org/abs/2407.21783}

\bibitem{esl}
T.~Hastie, R.~Tibshirani y J.~Friedman.
\newblock \emph{The Elements of Statistical Learning: Data Mining, Inference, and Prediction}, 2.ª ed.
\newblock Springer, 2009.
\newblock Caps.~2, 3, 4 y~7.
\newblock \url{https://hastie.su.domains/ElemStatLearn/}

\bibitem{hoerl}
A.~E. Hoerl y R.~W. Kennard.
\newblock Ridge regression: biased estimation for nonorthogonal problems.
\newblock \emph{Technometrics}, 12(1):55--67, 1970.

\bibitem{adamw}
I.~Loshchilov y F.~Hutter.
\newblock Decoupled weight decay regularization.
\newblock En \emph{International Conference on Learning Representations (ICLR)}, 2019.
\newblock \url{https://arxiv.org/abs/1711.05101}

\bibitem{mit6390}
MIT 6.390 (antes 6.036) Staff.
\newblock \emph{Introduction to Machine Learning}, notas del curso.
\newblock Massachusetts Institute of Technology. Cap.~2 «Regression» y cap.~4 «Classification».
\newblock \url{https://introml.mit.edu/notes/}

\bibitem{cs229}
A.~Ng y T.~Ma.
\newblock CS229 Lecture Notes (apuntes principales del curso).
\newblock Stanford University. Parte~I «Supervised learning» (caps.~1--3) y parte~III «Generalization and regularization» (caps.~8--9).
\newblock \url{https://cs229.stanford.edu/main_notes.pdf}

\bibitem{cs229ps}
Stanford CS229: Machine Learning.
\newblock \emph{Problem Set}~1.
\newblock Stanford University, verano de 2019.
\newblock \url{https://cs229.stanford.edu/summer2019/ps1.pdf}

\bibitem{strang}
G.~Strang.
\newblock \emph{Introduction to Linear Algebra}, 5.ª ed.
\newblock Wellesley-Cambridge Press, 2016.
\newblock Cap.~4 «Orthogonality».
\newblock \url{https://math.mit.edu/~gs/linearalgebra/}

\bibitem{mit1806}
G.~Strang.
\newblock 18.06SC: Linear Algebra.
\newblock Massachusetts Institute of Technology, MIT OpenCourseWare, otoño de 2011. Sesión~2.3 «Projection Matrices and Least Squares».
\newblock \url{https://ocw.mit.edu/courses/18-06sc-linear-algebra-fall-2011/}

\bibitem{lasso}
R.~Tibshirani.
\newblock Regression shrinkage and selection via the lasso.
\newblock \emph{Journal of the Royal Statistical Society, Series B}, 58(1):267--288, 1996.

\bibitem{llama2}
H.~Touvron \emph{et al.}
\newblock Llama 2: Open foundation and fine-tuned chat models.
\newblock arXiv:2307.09288, 2023.
\newblock \url{https://arxiv.org/abs/2307.09288}

\bibitem{d2l}
A.~Zhang, Z.~C. Lipton, M.~Li y A.~J. Smola.
\newblock \emph{Dive into Deep Learning}, versión 1.0.3.
\newblock Cambridge University Press, 2023.
\newblock §3.1 «Linear Regression», §3.6 «Generalization», §3.7 «Weight Decay» y §4.1 «Softmax Regression».
\newblock \url{https://d2l.ai/}

\end{thebibliography}


\end{document}
